\scanpage{19}
La décomposition envisagée (can., fonct.)
\[
\operatorname{DR}^1(A/S)\simeq\underline t_A^\vee+\underline t_{A^*}
\]
vaut pour tt préschéma abélien \uline{ordinaire} sur un $S$ de car. $p>0$, en $=$ temps que les expressions
\[
\underline t_A\simeq\unsure{\Phi|\mathcal L}_A\otimes_{\mathbb F_p}\underline{\mathrm O}_S,
\qquad \underline t_{A^*}\simeq\unsure{\Phi|\mathcal L}_{A^*}\otimes_{\mathbb F_p}\underline{\mathrm O}_S
\]
où $\unsure{\Phi|\mathcal L}_A$, $\unsure{\Phi|\mathcal L}_{A^*}$ sont des faisceaux \inserted{étales} \unsure{loc.} \deleted{const.} libres de rang $g$ ($=\dim\mathrm{rel}$ de $A/S$) de vectoriels sur $\mathbb F_p$. Si $S$ est le spectre d'un corps parfait,
\[
\begin{aligned}
\unsure{\Phi|\mathcal L}_{A^*}&=(\underline t_{A^*})^F
=H^1(\unsure{A|A^*},\underline{\mathrm O}_A)^F
=H^1(A,\mathbb Z/p\mathbb Z)\\
&=\operatorname{Hom}({}_pA,\mathbb Z/p\mathbb Z)
=\deleted{\unsure{\Phi}}\,{}_pA(k)^\vee
\end{aligned}
\]
\[
\left\{\begin{aligned}
\unsure{\Phi|\mathcal L}_{A^*}\deleted{\unsure{F}}(k)&={}_pA(k)^\vee\\
\unsure{\Phi|\mathcal L}_A(k)&={}_pA^*(k)^\vee
\end{aligned}\right.
\]
Donc
\[
\left\{\begin{aligned}
\underline t_A^\vee&=\unsure{\Phi|\mathcal L}_A\deleted{\unsure{(k)}}^\vee\otimes_{\mathbb F_p}k
={}_pA^*(k)\otimes_{\mathbb F_p}k\\
t_{A^*}&={}_pA(k)^\vee\otimes_{\mathbb F_p}k.
\end{aligned}\right.
\]
On a des relations analogues pour $S$ (de car $p>0$) quelconque: l'hyp. sur $A$ implique que
\diagram{\begin{tikzcd}{}_pA\arrow[r] & S\end{tikzcd}}
s'écrit de façon unique comme composé
\diagram{\begin{tikzcd}{}_pA\arrow[r] & ({}_pA)_{\text{ét}}\arrow[r] & S\end{tikzcd}}
de deux morphismes finis loc libres, le premier radiciel surj, le 2ième \uline{étale}.
\mspage{18}
